\frfilename{mt241.tex}
\versiondate{6.11.03}
\copyrightdate{1995}

\loadeusm

\def\chaptername{Ruang fungsi}
\def\sectionname{$\eusm L^0$ dan $L^0$}

\newsection{241}

Tujuan utama bab ini adalah membahas ruang-ruang $L^1$, $L^{\infty}$,
dan $L^p$ dalam tiga bagian berikutnya. Namun, akan lebih mudah jika semua
ruang tersebut dipandang sebagai subruang dari ruang yang lebih besar,
yaitu $L^0$, yang terdiri atas kelas-kelas ekuivalensi fungsi-fungsi yang
(secara maya) terukur; dalam bagian ini saya menghimpun fakta-fakta dasar
mengenai ruang linear terurut $L^0$.

Hampir dapat dikatakan bahwa prinsip pertama teori ukuran adalah bahwa
himpunan-himpunan berukuran nol sering kali dapat diabaikan; frasa
`himpunan terabaikan' itu sendiri menegaskan prinsip ini. Oleh karena itu,
dua fungsi yang sama hampir di mana-mana sering kali (tidak selalu!) dapat
diperlakukan sebagai identik. Ungkapan yang sesuai bagi gagasan ini adalah
membentuk ruang kelas-kelas ekuivalensi fungsi, dengan menyatakan bahwa dua
fungsi ekuivalen jika keduanya sama pada suatu himpunan koterabaikan. Inilah
landasan semua konstruksi dalam bab ini. Merupakan fakta yang mengagumkan
bahwa ruang-ruang kelas ekuivalensi yang dibangun dengan cara ini justru
lebih sesuai untuk masalah-masalah tertentu daripada ruang fungsi asalnya;
setelah tekniknya dikuasai, berpikir dalam ruang yang lebih abstrak itu pun
menjadi lebih mudah.


\leader{241A}{Ruang $\eusm L^0$: Definisi}
\cmmnt{ Kini tiba waktunya untuk memberi
nama kepada suatu himpunan fungsi yang telah digunakan lebih dari sekali.}
Misalkan $(X,\Sigma,\mu)$ suatu ruang ukur.
Saya menulis $\eusm L^0$, atau $\eusm L^0(\mu)$, untuk ruang
fungsi-fungsi bernilai riil $f$ yang didefinisikan pada subhimpunan
koterabaikan dari $X$ dan terukur secara maya\cmmnt{, artinya,
$f\restr E$ terukur untuk suatu himpunan koterabaikan $E\subseteq X$.
Ingat bahwa $f$ terukur secara maya terhadap $\mu$ jika dan hanya jika fungsi itu terukur terhadap
$\hat\Sigma$, dengan $\hat\Sigma$ kelengkapan dari $\Sigma$ (212Fa)}.

\leader{241B}{Sifat-sifat dasar}\cmmnt{ Jika $(X,\Sigma,\mu)$ suatu
ruang ukur sebarang, maka kita mempunyai fakta-fakta berikut, yang
bersesuaian dengan sifat-sifat dasar fungsi terukur yang tercantum dalam
\S121 Jilid 1. Saya membahasnya satu demi satu, sehingga jika Jilid 1 ada
di tangan Anda, Anda dapat melihat bagian mana yang harus dilewati.

\medskip

}{\bf (a)} Fungsi konstan bernilai riil yang didefinisikan hampir di
mana-mana dalam $X$ termasuk dalam $\eusm L^0$\prooflet{ (121Ea)}.

\spheader 241Bb $f+g\in\eusm L^0$ untuk semua $f$,
$g\in\eusm L^0$\prooflet{ (sebab jika $f\restr F$ dan $g\restr G$
terukur, maka $(f+g)\restr(F\cap G)=(f\restr F)+(g\restr G)$ terukur)
(121Eb)}.

\spheader 241Bc $cf\in\eusm L^0$ untuk semua $f\in\eusm L^0$,
$c\in\Bbb R$\prooflet{ (121Ec)}.

\spheader 241Bd $f\times g\in\eusm L^0$ untuk semua $f$,
$g\in\eusm L^0$\prooflet{ (121Ed)}.

\spheader 241Be Jika $f\in\eusm L^0$ dan $h:\Bbb R\to\Bbb R$ terukur
Borel, maka $hf\in\eusm L^0$\prooflet{ (121Eg)}.

\spheader 241Bf Jika $\sequencen{f_n}$ suatu barisan dalam $\eusm L^0$
dan $f=\lim_{n\to\infty}f_n$ didefinisikan\cmmnt{ (sebagai fungsi
bernilai riil)} hampir di mana-mana dalam $X$, maka
$f\in\eusm L^0$\prooflet{ (121Fa)}.

\spheader 241Bg Jika $\sequencen{f_n}$ suatu barisan dalam $\eusm L^0$
dan $f=\sup_{n\in\Bbb N}f_n$ didefinisikan\cmmnt{ (sebagai fungsi
bernilai riil)} hampir di mana-mana dalam $X$, maka
$f\in\eusm L^0$\prooflet{ (121Fb)}.

\spheader 241Bh Jika $\sequencen{f_n}$ suatu barisan dalam $\eusm L^0$
dan $f=\inf_{n\in\Bbb N}f_n$ didefinisikan\cmmnt{ (sebagai fungsi
bernilai riil)} hampir di mana-mana dalam $X$, maka
$f\in\eusm L^0$\prooflet{ (121Fc)}.

\spheader 241Bi Jika $\sequencen{f_n}$ suatu barisan dalam $\eusm L^0$
dan $f=\limsup_{n\to\infty}f_n$ didefinisikan\cmmnt{ (sebagai fungsi
bernilai riil)} hampir di mana-mana dalam $X$, maka
$f\in\eusm L^0$\prooflet{ (121Fd)}.

\spheader 241Bj Jika $\sequencen{f_n}$ suatu barisan dalam $\eusm L^0$
dan $f=\liminf_{n\to\infty}f_n$ didefinisikan\cmmnt{ (sebagai fungsi
bernilai riil)} hampir di mana-mana dalam $X$, maka
$f\in\eusm L^0$\prooflet{ (121Fe)}.

\spheader 241Bk $\eusm L^0$ tepat merupakan himpunan fungsi-fungsi
bernilai riil, yang didefinisikan pada subhimpunan-subhimpunan $X$, dan
sama hampir di mana-mana dengan suatu fungsi terukur terhadap $\Sigma$ dari
$X$ ke $\Bbb R$.
\prooflet{\Prf\ (i) Jika $g:X\to\Bbb R$ terukur terhadap $\Sigma$ dan
$f\eae g$, maka $F=\{x:x\in\dom f,\,f(x)=g(x)\}$ koterabaikan dan
$f\restr F=g\restr F$ terukur (121Eh), sehingga $f\in\eusm L^0$. (ii)
Jika $f\in\eusm L^0$, ambil himpunan koterabaikan $E\subseteq X$
sedemikian sehingga $f\restr E$ terukur. Maka $D=E\cap\dom f$
koterabaikan dan $f\restr D$ terukur, sehingga terdapat fungsi terukur
$h:X\to\Bbb R$ yang sama dengan $f$ pada $D$ (121I); dan $h\eae f$.
\ \Qed}

\leader{241C}{Ruang $L^0$: Definisi} Misalkan $(X,\Sigma,\mu)$ suatu
ruang ukur sebarang. Maka $\eae$ merupakan relasi ekuivalensi pada
$\eusm L^0$. Tuliskan $L^0$, atau $L^0(\mu)$, untuk himpunan kelas-kelas
ekuivalensi dalam $\eusm L^0$ terhadap $\eae$. Untuk $f\in\eusm L^0$,
tuliskan $f^{\ssbullet}$ untuk kelas ekuivalensinya dalam $L^0$.

\leader{241D}{Struktur linear $L^0$} Misalkan $(X,\Sigma,\mu)$ suatu
ruang ukur sebarang, dan tetapkan $\eusm L^0=\eusm L^0(\mu)$,
$L^0=L^0(\mu)$.

\spheader 241Da Jika $f_1$, $f_2$, $g_1$, $g_2\in\eusm L^0$,
$f_1\eae f_2$ dan $g_1\eae g_2$, maka $f_1+g_1\eae f_2+g_2$. Oleh
karena itu, kita dapat mendefinisikan penjumlahan pada $L^0$ dengan
menetapkan $f^{\ssbullet}+g^{\ssbullet}=(f+g)^{\ssbullet}$ untuk semua
$f$, $g\in\eusm L^0$.

\spheader 241Db Jika $f_1$, $f_2\in\eusm L^0$ dan $f_1\eae f_2$,
maka $cf_1\eae cf_2$ untuk setiap $c\in\Bbb R$. Oleh karena itu, kita
dapat mendefinisikan perkalian skalar pada $L^0$ dengan menetapkan
$c\cdot f^{\ssbullet}=(cf)^{\ssbullet}$ untuk semua $f\in\eusm L^0$
dan $c\in\Bbb R$.


\spheader 241Dc\cmmnt{ Kini} $L^0$ merupakan ruang linear atas
$\Bbb R$, dengan nol $\tbf{0}^{\ssbullet}$, dengan $\tbf{0}$ fungsi
berdomain $X$ dan bernilai konstan $0$, serta negatif
$-(f^{\ssbullet})=(-f)^{\ssbullet}$.
\prooflet{\Prf\ (i)

\Centerline{$f+(g+h)=(f+g)+h$ untuk semua $f$, $g$, $h\in\eusm L^0$,}

\noindent sehingga


\Centerline{$u+(v+w)=(u+v)+w$ untuk semua $u$, $v$, $w\in L^0$.}

(ii)

\Centerline{$f+\tbf{0}=\tbf{0}+f=f$ untuk setiap $f\in\eusm L^0$,}

\noindent sehingga

\Centerline{$u+\tbf{0}^{\ssbullet}=\tbf{0}^{\ssbullet}+u=u$ untuk
setiap $u\in L^0$.}

(iii)

\Centerline{$f+(-f)\eae\tbf{0}$ untuk setiap $f\in\eusm L^0$,}

\noindent sehingga
\Centerline{$f^{\ssbullet}+(-f)^{\ssbullet}=\tbf{0}^{\ssbullet}$ untuk
setiap $f\in\eusm L^0$.}

(iv)

\Centerline{$f+g=g+f$ untuk semua $f$, $g\in\eusm L^0$,}

\noindent sehingga

\Centerline{$u+v=v+u$ untuk semua $u$,
$v\in L^0$.}

(v)

\Centerline{$c(f+g)=cf+cg$ untuk semua $f$, $g\in\eusm L^0$ dan
$c\in\Bbb R$,}

\noindent sehingga

\Centerline{$c(u+v)=cu+cv$ untuk semua $u$, $v\in L^0$ dan $c\in\Bbb
R$.}

(vi)

\Centerline{$(a+b)f=af+bf$ untuk semua $f\in\eusm L^0$ dan $a$,
$b\in\Bbb R$,}

\noindent sehingga

\Centerline{$(a+b)u=au+bu$ untuk semua $u\in L^0$ dan $a$, $b\in\Bbb R$.}

(vii)

\Centerline{$(ab)f=a(bf)$ untuk semua $f\in\eusm L^0$ dan $a$,
$b\in\Bbb R$,}


\noindent sehingga


\Centerline{$(ab)u=a(bu)$ untuk semua $u\in L^0$ dan $a$, $b\in\Bbb R$.}

(viii)

\Centerline{$1f=f$ untuk semua $f\in\eusm L^0$,}

\noindent sehingga

\Centerline{$1u=u$ untuk semua $u\in L^0$. \Qed}
}%end of prooflet

\leader{241E}{Struktur urutan $L^0$} Misalkan $(X,\Sigma,\mu)$ suatu
ruang ukur sebarang dan tetapkan $\eusm L^0=\eusm L^0(\mu)$,
$L^0=L^0(\mu)$.

\spheader 241Ea Jika $f_1$, $f_2$, $g_1$, $g_2\in\eusm L^0$,
$f_1\eae f_2$, $g_1\eae g_2$, dan $f_1\leae g_1$, maka
$f_2\leae g_2$. Oleh karena itu, kita dapat mendefinisikan relasi $\le$
pada $L^0$ dengan mengatakan bahwa $f^{\ssbullet}\le g^{\ssbullet}$ jika dan hanya jika
$f\leae g$.

\spheader 241Eb\cmmnt{ Kini} $\le$ merupakan urutan parsial pada $L^0$.
\prooflet{\Prf\ (i) Jika $f$, $g$, $h\in\eusm L^0$, $f\leae g$, dan
$g\leae h$, maka $f\leae h$. Oleh karena itu, $u\le w$ kapan pun $u$,
$v$, $w\in L^0$, $u\le v$, dan $v\le w$. (ii) Jika $f\in\eusm L^0$,
maka $f\leae f$; jadi $u\le u$ untuk setiap $u\in L^0$. (iii) Jika
$f$, $g\in\eusm L^0$, $f\leae g$, dan $g\leae f$, maka $f\eae g$,
sehingga jika $u\le v$ dan $v\le u$, maka $u=v$.\ \Qed}

\spheader 241Ec\cmmnt{ Bahkan} $L^0$, dengan $\le$, merupakan
{\bf ruang linear terurut parsial}, yakni ruang linear (riil) dengan
urutan parsial $\le$ sedemikian sehingga

\qquad jika $u\le v$, maka $u+w\le v+w$ untuk setiap $w$,

\qquad jika $0\le u$, maka $0\le cu$ untuk setiap $c\ge 0$.

\noindent\prooflet{\Prf\ (i) Jika $f$, $g$, $h\in\eusm L^0$ dan
$f\leae g$, maka $f+h\leae g+h$. (ii) Jika $f\in\eusm L^0$ dan
$f\ge 0$ hampir di mana-mana, maka $cf\ge 0$ hampir di mana-mana untuk
setiap $c\ge 0$.\ \Qed}

\spheader 241Ed\cmmnt{ Lebih jauh lagi:} $L^0$ merupakan {\bf ruang
Riesz} atau {\bf kisi vektor}, yakni ruang linear terurut parsial
sedemikian sehingga $u\vee v=\sup\{u,v\}$ dan
$u\wedge v=\inf\{u,v\}$ terdefinisi untuk semua $u$, $v\in L^0$.
\prooflet{\Prf\ Ambil $f$, $g\in\eusm L^0$ sedemikian sehingga
$f^{\ssbullet}=u$ dan $g^{\ssbullet}=v$. Maka $f\vee g$,
$f\wedge g\in\eusm L^0$, dengan menuliskan

\Centerline{$(f\vee g)(x)=\max(f(x),g(x))$, \quad $(f\wedge
g)(x)=\min(f(x),g(x))$}

\noindent untuk $x\in\dom f\cap\dom g$. (Bandingkan 241Bg-h.) Kini,
untuk sebarang $h\in\eusm L^0$, kita mempunyai

\Centerline{$f\vee g\leae h\iff f\leae h$ dan $g\leae h$,}

\Centerline{$h\leae f\wedge g\iff h\leae f$ dan $h\leae g$,}

\noindent sehingga untuk sebarang $w\in L^0$ kita mempunyai

\Centerline{$(f\vee g)^{\ssbullet}\le w\iff u\le w$ dan $v\le w$,}

\Centerline{$w\le(f\wedge g)^{\ssbullet}\iff w\le u$ dan $w\le v$.}

\noindent Jadi, kita memperoleh

\Centerline{$(f\vee g)^{\ssbullet}=\sup\{u,v\}=u\vee v$,
\quad$(f\wedge g)^{\ssbullet}=\inf\{u,v\}=u\wedge v$}

\noindent dalam $L^0$.\ \Qed}

\spheader 241Ee Khususnya, untuk sebarang $u\in L^0$ kita dapat
membicarakan $|u|=u\vee(-u)$; jika $f\in\eusm L^0$, maka
$|f^{\ssbullet}|=|f|^{\ssbullet}$.

Jika\cmmnt{ $f$, $g\in\eusm L^0$,} $c\in\Bbb R$, maka

\prooflet{
\Centerline{$|cf|=|c||f|$, \quad$f\vee g=\Bover12(f+g+|f-g|)$,}

\Centerline{
$f\wedge g=\Bover12(f+g-|f-g|)$, \quad$|f+g|\leae|f|+|g|$,}

\noindent sehingga
}

\Centerline{$|cu|=|c||u|$, \quad$u\vee v=\Bover12(u+v+|u-v|)$,}

\Centerline{$u\wedge v=\Bover12(u+v-|u-v|)$, \quad$|u+v|\le|u|+|v|$}

\noindent untuk semua $u$, $v\in L^0$.

\spheader 241Ef\cmmnt{ Suatu notasi khusus sering kali berguna.}
Jika $f$ fungsi bernilai riil, tetapkan $f^+(x)=\max(f(x),0)$,
$f^-(x)=\max(-f(x),0)$ untuk $x\in\dom f$, sehingga

\Centerline{$f=f^+-f^-$, \quad$|f|=f^++f^-=f^+\vee f^-$,}

\noindent dengan semua fungsi tersebut didefinisikan pada $\dom f$.
Dalam $L^0$, operasi-operasi yang bersesuaian adalah $u^+=u\vee 0$,
$u^-=(-u)\vee 0$, dan kita mempunyai

\Centerline{$u=u^+-u^-$,
\quad$|u|=u^++u^-=u^+\vee u^-$,
\quad$u^+\wedge u^-=0$.}

\spheader 241Eg\dvro{ Jika}{ Barangkali hal berikut sudah jelas, tetapi
tetap saya nyatakan: jika} $u\ge 0$ dalam $L^0$, maka terdapat
$f\ge 0$ dalam $\eusm L^0$ sedemikian sehingga $f^{\ssbullet}=u$.
\prooflet{\Prf\ Ambil sebarang $g\in\eusm L^0$ sedemikian sehingga
$u=g^{\ssbullet}$, dan tetapkan $f=g\vee\tbf{0}$.\ \Qed}


\leader{241F}{Ruang Riesz}\cmmnt{ Terdapat teori abstrak yang luas
mengenai ruang Riesz, yang menurut saya sebaiknya kita kesampingkan untuk
saat ini; uraian umum dapat ditemukan dalam {\smc Luxemburg \& Zaanen 71}
dan {\smc Zaanen 83}; buku saya sendiri {\smc Fremlin 74} membahas bahan
elementernya, dan Bab 35 dalam jilid berikutnya mengulangi gagasan-gagasan
yang paling esensial. Untuk keperluan kita di sini, kita hanya memerlukan
beberapa definisi dan sejumlah hasil sederhana yang paling mudah dibuktikan
untuk kasus-kasus khusus yang kita butuhkan, tanpa mengacu pada teori umum.

\medskip

} {\bf (a)} Suatu ruang Riesz $U$ disebut {\bf Archimedes} jika kapan
pun $u\in U$, $u>0$\cmmnt{ (artinya, $u\ge 0$ dan $u\ne 0$),} dan
$v\in U$, terdapat $n\in\Bbb N$ sedemikian sehingga $nu\not\le v$.

\spheader 241Fb Suatu ruang Riesz $U$ disebut {\bf lengkap-$\sigma$
Dedekind}\cmmnt{ (atau {\bf lengkap-$\sigma$ menurut urutan}, atau
{\bf lengkap-$\sigma$})} jika setiap himpunan tak kosong terhitung
$A\subseteq U$ yang terbatas di atas mempunyai batas atas terkecil dalam
$U$.

\spheader 241Fc Suatu ruang Riesz $U$ disebut {\bf lengkap Dedekind}
(atau {\bf lengkap menurut urutan}, atau {\bf lengkap}) jika setiap
himpunan tak kosong $A\subseteq U$ yang terbatas di atas dalam $U$
mempunyai batas atas terkecil dalam $U$.

\leader{241G}{}\cmmnt{ Kini kita mempunyai sifat-sifat penting $L^0$
berikut.

\medskip

\noindent}{\bf Teorema} Misalkan $(X,\Sigma,\mu)$ suatu ruang ukur.
Tetapkan $L^0=L^0(\mu)$.

(a) $L^0$ bersifat Archimedes dan lengkap-$\sigma$ Dedekind.

(b) Jika $(X,\Sigma,\mu)$ semihingga, maka $L^0$ lengkap Dedekind jika dan hanya jika
$(X,\Sigma,\mu)$ dapat dilokalkan.

\proof{ Tetapkan $\eusm L^0=\eusm L^0(\mu)$.

\medskip

{\bf (a)(i)} Jika $u$, $v\in L^0$ dan $u>0$, nyatakan $u$ sebagai
$f^{\ssbullet}$ dan $v$ sebagai $g^{\ssbullet}$ dengan $f$,
$g\in\eusm L^0$. Maka $E=\{x:x\in\dom f,\,f(x)>0\}$ tidak terabaikan.
Jadi, terdapat $n\in\Bbb N$ sedemikian sehingga

\Centerline{$E_n=\{x:x\in\dom f\cap\dom g,\,nf(x)>g(x)\}$}

\noindent tidak terabaikan, sebab
$E\cap\dom g\subseteq\bigcup_{n\in\Bbb N}E_n$. Tetapi sekarang
$nu\not\le v$. Karena $u$ dan $v$ sebarang, $L^0$ bersifat Archimedes.

\medskip

\quad{\bf (ii)} Sekarang misalkan $A\subseteq L^0$ suatu himpunan tak
kosong terhitung dengan batas atas $w$ dalam $L^0$. Nyatakan $A$ sebagai
$\{f_n^{\ssbullet}:n\in\Bbb N\}$ dengan $\sequencen{f_n}$ suatu barisan
dalam $\eusm L^0$, dan $w$ sebagai $h^{\ssbullet}$ dengan
$h\in\eusm L^0$. Tetapkan $f=\sup_{n\in\Bbb N}f_n$. Maka $f(x)$
terdefinisi dalam $\Bbb R$ di setiap titik
$x\in\dom h\cap\bigcap_{n\in\Bbb N}\dom f_n$ sedemikian sehingga
$f_n(x)\le h(x)$ untuk setiap $n\in\Bbb N$, yakni untuk hampir setiap
$x\in X$; jadi $f\in\eusm L^0$ (241Bg). Tetapkan
$u=f^{\ssbullet}\in L^0$. Jika $v\in L^0$, katakanlah
$v=g^{\ssbullet}$ dengan $g\in\eusm L^0$, maka

$$\eqalign{u_n\le v&\text{ untuk setiap }n\in\Bbb N\cr
&\iff\text{ untuk setiap }n\in\Bbb N,\,f_n\leae g\cr
&\iff\text{ untuk hampir setiap }x\in X,\,f_n(x)\le g(x)
      \text{ untuk setiap }n\in\Bbb N\cr
&\iff f\leae g
\iff u\le v.\cr}$$

\noindent Jadi, $u=\sup_{n\in\Bbb N}u_n$ dalam $L^0$. Karena $A$
sebarang, $L^0$ lengkap-$\sigma$ Dedekind.

\medskip

{\bf (b)(i)} Misalkan $(X,\Sigma,\mu)$ dapat dilokalkan. Misalkan
$A\subseteq L^0$ sebarang himpunan tak kosong dengan batas atas
$w_0\in L^0$. Tetapkan

\Centerline{$\Cal A=\{f:f$ fungsi terukur dari $X$ ke
$\Bbb R,\,f^{\ssbullet}\in A\}$;}

\noindent maka setiap anggota $A$ berbentuk $f^{\ssbullet}$ untuk suatu
$f\in\Cal A$ (241Bk). Untuk setiap $q\in\Bbb Q$, misalkan $\Cal E_q$
keluarga subhimpunan $X$ yang dapat dinyatakan dalam bentuk
$\{x:f(x)\ge q\}$ untuk suatu $f\in\Cal A$; maka
$\Cal E_q\subseteq\Sigma$. Karena $(X,\Sigma,\mu)$ dapat dilokalkan,
terdapat himpunan $F_q\in\Sigma$ yang merupakan supremum esensial bagi
$\Cal E_q$. Untuk $x\in X$, tetapkan

\Centerline{$g^*(x)=\sup\{q:q\in\Bbb Q,\,x\in F_q\}$,}

\noindent dengan mengizinkan $\infty$ sebagai supremum suatu himpunan
yang tidak terbatas di atas, dan $-\infty$ sebagai $\sup\emptyset$.
Maka

\Centerline{$\{x:g^*(x)>a\}=\bigcup_{q\in\Bbb Q,q>a}F_q\in\Sigma$}

\noindent untuk setiap $a\in\Bbb R$.

Jika $f\in\Cal A$, maka $f\leae g^*$. \Prf\ Untuk setiap
$q\in\Bbb Q$, tetapkan

\Centerline{$E_q=\{x:f(x)\ge q\}\in\Cal E_q$;}

\noindent maka $E_q\setminus F_q$ terabaikan. Tetapkan
$H=\bigcup_{q\in\Bbb Q}(E_q\setminus F_q)$. Jika $x\in X\setminus H$,
maka

\Centerline{$f(x)\ge q\Longrightarrow g^*(x)\ge q$,}

\noindent sehingga $f(x)\le g^*(x)$; jadi $f\leae g^*$.\ \Qed

Jika $h:X\to\Bbb R$ terukur dan $u\le h^{\ssbullet}$ untuk setiap
$u\in A$, maka $g^*\leae h$. \Prf\ Tetapkan
$G_q=\{x:h(x)\ge q\}$ untuk setiap $q\in\Bbb Q$. Jika
$E\in\Cal E_q$, terdapat $f\in\Cal A$ sedemikian sehingga
$E=\{x:f(x)\ge q\}$; kini $f\leae h$, sehingga
$E\setminus G_q\subseteq\{x:f(x)>h(x)\}$ terabaikan. Karena $F_q$
merupakan supremum esensial bagi $\Cal E_q$, $F_q\setminus G_q$
terabaikan; dan ini berlaku untuk setiap $q\in\Bbb Q$. Akibatnya,

\Centerline{$\{x:h(x)<g^*(x)\}
\subseteq\bigcup_{q\in\Bbb Q}F_q\setminus G_q$}

\noindent terabaikan, dan $g^*\leae h$.\ \Qed

Kini ingat bahwa kita mengandaikan $A\ne\emptyset$ dan bahwa $A$
mempunyai batas atas $w_0\in L^0$. Ambil sebarang $f_0\in\Cal A$ dan
fungsi terukur $h_0:X\to\Bbb R$ sedemikian sehingga
$h_0^{\ssbullet}=w_0$; maka $f\leae h_0$ untuk setiap $f\in\Cal A$,
sehingga $f_0\leae g^*\leae h_0$, dan $g^*$ harus berhingga hampir di
mana-mana. Dengan menetapkan $g(x)=g^*(x)$ ketika $g^*(x)\in\Bbb R$,
kita memperoleh $g\in\eusm L^0$ dan $g\eae g^*$, sehingga

\Centerline{$f\leae g\leae h$}

\noindent kapan pun $f$, $h$ merupakan fungsi terukur dari $X$ ke
$\Bbb R$, $f^{\ssbullet}\in A$, dan $h^{\ssbullet}$ merupakan batas
atas bagi $A$; yakni,

\Centerline{$u\le g^{\ssbullet}\le w$}

\noindent kapan pun $u\in A$ dan $w$ merupakan batas atas bagi $A$.
Tetapi ini berarti $g^{\ssbullet}$ merupakan batas atas terkecil bagi
$A$ dalam $L^0$. Karena $A$ sebarang, $L^0$ lengkap Dedekind.

\medskip

{\bf (ii)} Misalkan $L^0$ lengkap Dedekind. Kita mengandaikan bahwa
$(X,\Sigma,\mu)$ semihingga. Misalkan $\Cal E$ sebarang subhimpunan
dari $\Sigma$. Tetapkan

\Centerline{$A=\{0\}\cup\{(\chi E)^{\ssbullet}:E\in\Cal E\}\subseteq
L^0$.}

\noindent Maka $A$ terbatas di atas oleh $(\chi X)^{\ssbullet}$,
sehingga mempunyai batas atas terkecil $w\in L^0$. Nyatakan $w$ sebagai
$h^{\ssbullet}$ dengan $h:X\to\Bbb R$ terukur, dan tetapkan
$F=\{x:h(x)>0\}$. Maka $F$ merupakan supremum esensial bagi $\Cal E$
dalam $\Sigma$. \Prf\ {\bf (${\alpha}$)} Jika $E\in\Cal E$, maka
$(\chi E)^{\ssbullet}\le w$, sehingga $\chi E\leae h$, yakni
$h(x)\ge 1$ untuk hampir setiap $x\in E$, dan
$E\setminus F\subseteq\{x:x\in E,\,h(x)<1\}$ terabaikan.
{\bf (${\beta}$)} Jika $G\in\Sigma$ dan $E\setminus G$ terabaikan untuk
setiap $E\in\Cal E$, maka $\chi E\leae\chi G$ untuk setiap
$E\in\Cal E$, yakni $(\chi E)^{\ssbullet}\le(\chi G)^{\ssbullet}$
untuk setiap $E\in\Cal E$; jadi $w\le(\chi G)^{\ssbullet}$, yakni
$h\leae\chi G$. Oleh karena itu,
$F\setminus G\subseteq\{x:h(x)>(\chi G)(x)\}$ terabaikan.\ \Qed

Karena $\Cal E$ sebarang, $(X,\Sigma,\mu)$ dapat dilokalkan.
}%end of proof of 241G

\leader{241H}{Struktur perkalian $L^0$} Misalkan $(X,\Sigma,\mu)$ suatu
ruang ukur sebarang; tuliskan $L^0=L^0(\mu)$,
$\eusm L^0=\eusm L^0(\mu)$.

\spheader 241Ha Jika $f_1$, $f_2$, $g_1$, $g_2\in\eusm L^0$,
$f_1\eae f_2$, dan $g_1\eae g_2$, maka
$f_1\times g_1\eae f_2\times g_2$. Oleh karena itu, kita dapat
mendefinisikan perkalian pada $L^0$ dengan menetapkan
$f^{\ssbullet}\times g^{\ssbullet}=(f\times g)^{\ssbullet}$ untuk semua
$f$, $g\in\eusm L^0$.

\spheader 241Hb\dvro{ Untuk}{ Kini mudah diperiksa bahwa, untuk} semua
$u$, $v$, $w\in L^0$ dan $c\in\Bbb R$,

\qquad $u\times(v\times w)=(u\times v)\times w$,

\qquad $u\times e=e\times u=u$,

\noindent dengan $e=\chi X^{\ssbullet}$ kelas ekuivalensi fungsi yang
bernilai konstan $1$,

\qquad $c(u\times v)=cu\times v=u\times cv$,

\qquad $u\times(v+w)=(u\times v)+(u\times w)$,

\qquad $(u+v)\times w=(u\times w)+(v\times w)$,

\qquad $u\times v=v\times u$,

\qquad $|u\times v|=|u|\times |v|$,

\qquad $u\times v=0$ jika dan hanya jika $|u|\wedge|v|=0$,

\qquad $|u|\le|v|$ jika dan hanya jika terdapat $w$ sedemikian sehingga
$|w|\le e$ dan $u=v\times w$.

\leader{241I}{Aksi fungsi Borel pada $L^0$} Misalkan
$(X,\Sigma,\mu)$ suatu ruang ukur dan $h:\Bbb R\to\Bbb R$ suatu fungsi
terukur Borel. Maka\cmmnt{ $hf\in\eusm L^0=\eusm L^0(\mu)$ untuk
setiap $f\in\eusm L^0$\prooflet{ (241Be)} dan $hf\eae hg$ kapan pun
$f\eae g$. Jadi} kita mempunyai fungsi $\bar h:L^0\to L^0$ yang
didefinisikan dengan menetapkan
$\bar h(f^{\ssbullet})=(hf)^{\ssbullet}$ untuk setiap
$f\in\eusm L^0$. \cmmnt{Sebagai contoh, jika $u\in L^0$ dan $p\ge 1$,
kita dapat meninjau $|u|^p=\bar h(u)$ dengan $h(x)=|x|^p$ untuk
$x\in\Bbb R$.}



\leader{241J}{$L^0$ kompleks}\cmmnt{ Gagasan-gagasan dalam bab ini,
seperti halnya dalam Bab 22-23, sering diterapkan pada ruang-ruang yang
berdasarkan fungsi bernilai kompleks, bukan fungsi bernilai riil.}
Misalkan $(X,\Sigma,\mu)$ suatu ruang ukur.

\spheader 241Ja Kita dapat menuliskan
$\eusm L^0_{\Bbb C}=\eusm L^0_{\Bbb C}(\mu)$ untuk ruang fungsi-fungsi
bernilai kompleks $f$ sedemikian sehingga $\dom f$ merupakan subhimpunan
koterabaikan dari $X$ dan terdapat subhimpunan koterabaikan
$E\subseteq X$ sedemikian sehingga $f\restr E$ terukur; yakni,
sedemikian sehingga bagian riil dan bagian imajiner $f$ keduanya termasuk
dalam $\eusm L^0(\mu)$. \cmmnt{Selanjutnya,}
$L^0_{\Bbb C}=L^0_{\Bbb C}(\mu)$ merupakan ruang kelas-kelas
ekuivalensi dalam $\eusm L^0_{\Bbb C}$ terhadap relasi ekuivalensi
$\eae$.

\spheader 241Jb Dengan menggunakan rumus yang sama seperti dalam 241D,
mudah untuk mendeskripsikan penjumlahan dan perkalian skalar yang menjadikan
$L^0_{\Bbb C}$ ruang linear atas $\Bbb C$. Kita\cmmnt{ tidak lagi
mempunyai struktur urutan yang persis sama, tetapi kita} dapat
mengidentifikasi suatu `bagian riil', yaitu

\Centerline{$\{f^{\ssbullet}:f\in\eusm L^0_{\Bbb C}$ bernilai riil
hampir di mana-mana$\}$,}

\noindent yang \cmmnt{jelas }dapat diidentifikasi dengan ruang linear
riil $L^0$, serta peta-peta yang bersesuaian
$u\mapsto\Real(u)$, $u\mapsto\Imag(u):L^0_{\Bbb C}\to L^0$ sedemikian
sehingga $u=\Real(u)+i\Imag(u)$ untuk setiap $u$. Selain itu, kita
mempunyai pengertian `modulus', dengan menuliskan

\Centerline{$|f^{\ssbullet}|=|f|^{\ssbullet}\in L^0$ untuk setiap
$f\in\eusm L^0_{\Bbb C}$,}

\noindent yang memenuhi relasi dasar $|cu|=|c||u|$,
$|u+v|\le|u|+|v|$ untuk $u$, $v\in L^0_{\Bbb C}$ dan
$c\in\Bbb C$\cmmnt{, seperti dalam 241Ef}.
Kita tentu\cmmnt{ saja} masih mempunyai perkalian pada $L^0_{\Bbb C}$,
dan semua rumus dalam 241H tetap berlaku untuknya.

\spheader 241Jc\cmmnt{ Fakta berikut berguna.} Untuk sebarang
$u\in L^0_{\Bbb C}$, $|u|$ merupakan supremum dalam $L^0$ dari
$\{\Real(\zeta u):\zeta\in\Bbb C,\,|\zeta|=1\}$. \prooflet{\Prf\
(i) Jika $|\zeta|=1$, maka $\Real(\zeta u)\le|\zeta u|=|u|$. Jadi,
$|u|$ merupakan batas atas dari
$\{\Real(\zeta u):|\zeta|=1\}$. (ii) Jika $v\in L^0$ dan
$\Real(\zeta u)\le v$ kapan pun $|\zeta|=1$, nyatakan $u$, $v$ sebagai
$f^{\ssbullet}$, $g^{\ssbullet}$ dengan $f:X\to\Bbb C$ dan
$g:X\to\Bbb R$ terukur. Untuk sebarang $q\in\Bbb Q$, $x\in X$,
tetapkan $f_q(x)=\Real(e^{iq}f(x))$. Maka $f_q\leae g$. Oleh karena itu,
$H=\{x:f_q(x)\le g(x)$ untuk setiap $q\in\Bbb Q\}$ koterabaikan.
Tetapi tentu saja $H=\{x:|f(x)|\le g(x)\}$, sehingga $|f|\leae g$ dan
$|u|\le v$. Karena $v$ sebarang, $|u|$ merupakan batas atas terkecil
dari $\{\Real(\zeta u):|\zeta|=1\}$.\ \Qed}

\exercises{
\leader{241X}{Latihan dasar $\pmb{>}$(a)}
%\spheader 241Xa
Misalkan $X$ suatu himpunan, dan misalkan $\mu$ ukuran cacah pada $X$
(112Bd). Tunjukkan bahwa $L^0(\mu)$ dapat diidentifikasi dengan
$\eusm L^0(\mu)=\Bbb R^X$.
%241C

\sqheader 241Xb Misalkan $(X,\Sigma,\mu)$ suatu ruang ukur dan
$\hat\mu$ kelengkapan dari $\mu$. Tunjukkan bahwa
$\eusm L^0(\mu)=\eusm L^0(\hat\mu)$ dan
$L^0(\mu)=L^0(\hat\mu)$.
%241C

\spheader 241Xc Misalkan $(X,\Sigma,\mu)$ suatu ruang ukur. (i)
Tunjukkan bahwa untuk setiap $u\in L^0(\mu)$ kita dapat mendefinisikan
suatu ukuran luar $\theta_u:\Cal P\Bbb R\to[0,\infty]$ dengan menuliskan
$\theta_u(A)=\mu^*f^{-1}[A]$ kapan pun $A\subseteq\Bbb R$ dan
$f\in\eusm L^0(\mu)$ sedemikian sehingga $f^{\ssbullet}=u$. (ii)
Tunjukkan bahwa ukuran yang didefinisikan dari $\theta_u$ dengan metode
\Caratheodory\ mengukur setiap subhimpunan Borel dari $\Bbb R$.
%241C

\spheader 241Xd Misalkan
$\langle(X_i,\Sigma_i,\mu_i)\rangle_{i\in I}$ suatu keluarga ruang ukur,
dengan jumlah langsung $(X,\Sigma,\mu)$ (214L). (i) Dengan menuliskan
$\phi_i:X_i\to X$ untuk peta kanonik (dalam konstruksi 214L,
$\phi_i(x)=(x,i)$ untuk $x\in X_i$), tunjukkan bahwa
$f\mapsto\langle f\phi_i\rangle_{i\in I}$ merupakan bijeksi antara
$\eusm L^0(\mu)$ dan $\prod_{i\in I}\eusm L^0(\mu_i)$. (ii) Tunjukkan
bahwa ini bersesuaian dengan suatu bijeksi antara $L^0(\mu)$ dan
$\prod_{i\in I}L^0(\mu_i)$.
%241C

\spheader 241Xe Misalkan $U$ suatu ruang Riesz lengkap-$\sigma$
Dedekind dan $A\subseteq U$ suatu himpunan tak kosong terhitung yang
terbatas di bawah dalam $U$. Tunjukkan bahwa $\inf A$ terdefinisi dalam
$U$.
%241F

\spheader 241Xf Misalkan $U$ suatu ruang Riesz lengkap Dedekind dan
$A\subseteq U$ suatu himpunan tak kosong yang terbatas di bawah dalam
$U$. Tunjukkan bahwa $\inf A$ terdefinisi dalam $U$.
%241F

\spheader 241Xg Misalkan $(X,\Sigma,\mu)$ dan $(Y,\Tau,\nu)$
ruang-ruang ukur, dan $\phi:X\to Y$ suatu \imp\ fungsi. (i) Tunjukkan
bahwa kita mempunyai peta $T:L^0(\nu)\to L^0(\mu)$ yang didefinisikan
dengan menetapkan $Tg^{\ssbullet}=(g\phi)^{\ssbullet}$ untuk setiap
$g\in\eusm L^0(\nu)$. (ii) Tunjukkan bahwa $T$ linear, bahwa
$T(v\times w)=Tv\times Tw$ untuk semua $v$, $w\in L^0(\nu)$, dan bahwa
$T(\sup_{n\in\Bbb N}v_n)=\sup_{n\in\Bbb N}Tv_n$ kapan pun
$\sequencen{v_n}$ suatu barisan dalam $L^0(\nu)$ dengan batas atas dalam
$L^0(\nu)$.
%241H

\sqheader 241Xh Misalkan $(X,\Sigma,\mu)$ suatu ruang ukur. Andaikan
$r\ge 1$ dan $h:\BbbR^r\to\Bbb R$ suatu fungsi terukur Borel.
Tunjukkan bahwa terdapat fungsi $\bar h:L^0(\mu)^r\to L^0(\mu)$ yang
didefinisikan dengan menuliskan

\Centerline{$\bar h(f_1^{\ssbullet},\ldots,f_r^{\ssbullet})
=(h(f_1,\ldots,f_r))^{\ssbullet}$}

\noindent untuk $f_1,\ldots,f_r\in{\eusm L}^0(\mu)$.
%241I

\spheader 241Xi Misalkan $(X,\Sigma,\mu)$ suatu ruang ukur dan $g$,
$h$, $\sequencen{g_n}$ fungsi-fungsi terukur Borel dari $\Bbb R$ ke
dirinya sendiri; tuliskan $\bar g$, $\bar h$, $\bar g_n$ untuk
fungsi-fungsi yang bersesuaian dari $L^0=L^0(\mu)$ ke dirinya sendiri
(241I). (i) Tunjukkan bahwa

\Centerline{$\bar g(u)+\bar h(u)=\overline{g+h}(u)$,
\quad$\bar g(u)\times\bar h(u)=\overline{g\times h}(u)$,
\quad$\bar g(\bar h(u))=\overline{gh}(u)$}

\noindent untuk setiap $u\in L^0$. (ii) Tunjukkan bahwa jika
$g(t)\le h(t)$ untuk setiap $t\in\Bbb R$, maka
$\bar g(u)\le\bar h(u)$ untuk setiap $u\in L^0$. (iii) Tunjukkan bahwa
jika $g$ tak menurun, maka $\bar g(u)\le\bar g(v)$ kapan pun $u\le v$
dalam $L^0$. (iv) Tunjukkan bahwa jika
$h(t)=\sup_{n\in\Bbb N}g_n(t)$ untuk setiap $t\in\Bbb R$, maka
$\bar h(u)=\sup_{n\in\Bbb N}\bar g_n(u)$ dalam $L^0$ untuk setiap
$u\in L^0$.
%241I

\leader{241Y}{Latihan lanjutan (a)}
%\spheader 241Ya
Misalkan $U$ sebarang ruang Riesz. Untuk $u\in U$, tuliskan
$|u|=u\vee(-u)$, $u^+=u\vee 0$, $u^-=(-u)\vee 0$. Tunjukkan bahwa,
untuk sebarang $u$, $v\in U$,

\Centerline{$u=u^+-u^-$,\quad$|u|=u^++u^-=u^+\vee u^-$,
\quad$u^+\wedge u^-=0$,}

\Centerline{$u\vee v=\Bover12(u+v+|u-v|)=u+(v-u)^+$,}

\Centerline{$u\wedge v=\Bover12(u+v-|u-v|)=u-(u-v)^+$,}

\Centerline{$|u+v|\le|u|+|v|$.}
%241E

\spheader 241Yb Misalkan $U$ suatu ruang linear terurut parsial dan
$N$ suatu subruang linear dari $U$ sedemikian sehingga kapan pun $u$,
$u'\in N$ dan $u'\le v\le u$, maka $v\in N$. (i) Tunjukkan bahwa ruang
linear hasil bagi $U/N$ merupakan ruang linear terurut parsial jika kita
mengatakan bahwa $u^{\ssbullet}\le v^{\ssbullet}$ dalam $U/N$ jika dan hanya jika
terdapat $w\in N$ sedemikian sehingga $u\le v+w$ dalam $U$. (ii)
Tunjukkan bahwa dalam hal ini $U/N$ merupakan ruang Riesz jika $U$
merupakan ruang Riesz dan $|u|\in N$ untuk setiap $u\in N$.
%241E

\spheader 241Yc Misalkan $(X,\Sigma,\mu)$ suatu ruang ukur. Tuliskan
$\eusm L^0_{\Sigma}$ untuk ruang semua fungsi terukur dari $X$ ke
$\Bbb R$, dan $\eusm N$ untuk subruang dari $\eusm L^0_{\Sigma}$ yang
terdiri atas fungsi-fungsi terukur yang bernilai nol hampir di mana-mana.
(i) Tunjukkan bahwa $\eusm L^0_{\Sigma}$ merupakan ruang Riesz
lengkap-$\sigma$ Dedekind. (ii) Tunjukkan bahwa $L^0(\mu)$ dapat
diidentifikasi, sebagai ruang linear terurut, dengan hasil bagi
$\eusm L^0_{\Sigma}/\eusm N$ sebagaimana didefinisikan dalam 241Yb di
atas.
%241Yb 241E

\spheader 241Yd Tunjukkan bahwa setiap ruang Riesz lengkap-$\sigma$
Dedekind bersifat Archimedes.
%241F

\spheader 241Ye Suatu ruang Riesz $U$ dikatakan memiliki
{\bf sifat supremum terhitung} jika untuk setiap $A\subseteq U$ yang
mempunyai batas atas terkecil dalam $U$, terdapat $B\subseteq A$
terhitung sedemikian sehingga $\sup B=\sup A$. Tunjukkan bahwa jika
$(X,\Sigma,\mu)$ suatu ruang ukur semihingga, maka ruang itu
$\sigma$-hingga jika dan hanya jika $L^0(\mu)$ memiliki sifat supremum terhitung.
%241G

\spheader 241Yf
Misalkan $(X,\Sigma,\mu)$ suatu ruang ukur dan $\tilde\mu$ versi c.l.d.\
dari $\mu$ (213E). (i) Tunjukkan bahwa
$\eusm L^0(\mu)\subseteq\eusm L^0(\tilde\mu)$. (ii) Tunjukkan bahwa
inklusi ini mendefinisikan operator linear
$T:L^0(\mu)\to L^0(\tilde\mu)$ sedemikian sehingga
$T(u\times v)=Tu\times Tv$ untuk semua $u$, $v\in L^0(\mu)$. (iii)
Tunjukkan bahwa kapan pun $v>0$ dalam $L^0(\tilde\mu)$, terdapat
$u\ge 0$ dalam $L^0(\mu)$ sedemikian sehingga $0<Tu\le v$. (iv)
Tunjukkan bahwa $T(\sup A)=\sup T[A]$ kapan pun
$A\subseteq L^0(\mu)$ suatu himpunan tak kosong dengan batas atas terkecil
dalam $L^0(\mu)$. (v) Tunjukkan bahwa $T$ injektif jika dan hanya jika $\mu$
semihingga. (vi) Tunjukkan bahwa jika $\mu$ dapat dilokalkan, maka $T$
merupakan isomorfisme untuk struktur linear dan urutan dari $L^0(\mu)$
dan $L^0(\tilde\mu)$. \Hint{213Hb.}
%241H

\spheader 241Yg Misalkan $(X,\Sigma,\mu)$ suatu ruang ukur dan $Y$
sebarang subhimpunan dari $X$; misalkan $\mu_Y$ ukuran subruang pada $Y$.
(i) Tunjukkan bahwa
$\eusm L^0(\mu_Y)=\{f\restrp Y:f\in\eusm L^0(\mu)\}$. (ii) Tunjukkan
bahwa terdapat surjeksi kanonik $T:L^0(\mu)\to L^0(\mu_Y)$ yang
didefinisikan dengan menetapkan
$T(f^{\ssbullet})=(f\restrp Y)^{\ssbullet}$ untuk setiap
$f\in\eusm L^0(\mu)$, yang linear dan multiplikatif serta mempertahankan
supremum dan infimum hingga, sehingga (khususnya) $T(|u|)=|Tu|$ untuk
setiap $u\in L^0(\mu)$. (iii) Tunjukkan bahwa $T$ injektif jika dan hanya jika $Y$
mempunyai ukuran luar penuh.
%241H

\spheader 241Yh Andaikan, dalam 241Yg, bahwa $Y\in\Sigma$.
Jelaskan bagaimana $L^0(\mu_Y)$ dapat diidentifikasi (sebagai ruang linear
terurut) dengan subruang
$\{u:u\times\chi(X\setminus Y)^{\ssbullet}=0\}$ dari $L^0(\mu)$.
%241Yg 241H

\spheader 241Yi Misalkan $(X,\Sigma,\mu)$ suatu ruang ukur, dan
$h:\Bbb R\to\Bbb R$ suatu fungsi tak menurun yang kontinu dari kiri.
Tunjukkan bahwa jika $A\subseteq L^0=L^0(\mu)$ suatu himpunan tak kosong
dengan supremum $v\in L^0$, maka $\bar h(v)=\sup_{u\in A}\bar h(u)$,
dengan $\bar h:L^0\to L^0$ fungsi yang dideskripsikan dalam 241I.
%241I
}%end of exercises

\endnotes{
\Notesheader{241} Sebagaimana disiratkan dalam 241Ya dan 241Yd,
sifat-sifat elementer ruang $L^0$ yang mengisi sebagian besar bagian ini
saling bergantung dengan kuat; tidaklah sulit untuk mengembangkan teori
`aljabar Riesz' guna memadukan gagasan-gagasan dalam 241H dengan
gagasan-gagasan lainnya. (Sesungguhnya, saya menguraikan teori semacam itu
dalam \S352 di jilid berikutnya, dengan nama `aljabar-$f$'.)

Jika kita menuliskan $\eusm L^0_{\Sigma}$ untuk ruang fungsi-fungsi terukur
dari $X$ ke $\Bbb R$, maka $\eusm L^0_{\Sigma}$ juga merupakan ruang Riesz
lengkap-$\sigma$ Dedekind, dan $L^0$ dapat diidentifikasi dengan hasil bagi
$\eusm L^0_{\Sigma}/\eusm N$, dengan $\eusm N$ menyatakan himpunan
fungsi-fungsi dalam $\eusm L^0_{\Sigma}$ yang bernilai nol hampir di
mana-mana. (Untuk melakukannya dengan tepat, kita memerlukan teori hasil bagi
ruang linear terurut; lihat 241Yb-241Yc di atas.) Tentu saja $\eusm L^0$,
sebagaimana saya mendefinisikannya, tidak sepenuhnya merupakan ruang linear.
Saya memilih deskripsi $L^0$ yang sedikit lebih rumit sebagai ruang kelas
ekuivalensi dalam $\eusm L^0$, bukan dalam $\eusm L^0_{\Sigma}$, sebab
dalam praktik sering terjadi bahwa suatu anggota $L^0$ berasal dari anggota
$\eusm L^0$ yang entah tidak didefinisikan di setiap titik ruang dasarnya,
atau tidak sepenuhnya terukur; dan menyesuaikan fungsi semacam itu supaya
menjadi anggota $\eusm L^0_{\Sigma}$, sekalipun sepele, merupakan proses
sewenang-wenang yang menurut saya dapat menyimpangkan sifat sejati konstruksi
tersebut. Tentu saja argumen yang sama dapat digunakan untuk mendukung ruang
yang sedikit lebih besar, yaitu ruang $\eusm L^0_{\infty}$ dari fungsi
yang terukur secara maya terhadap $\mu$, bernilai $[-\infty,\infty]$, serta
didefinisikan dan berhingga hampir di mana-mana, dengan mengandalkan 135E,
bukan 121E-121F. Akan tetapi, saya berpendapat bahwa operasi membatasi suatu
fungsi dalam $\eusm L^0_{\infty}$ pada himpunan tempat fungsi itu berhingga
{\it bukan} sewenang-wenang, melainkan kanonik dan sepenuhnya alami.

Ketika membaca uraian di atas -- atau, dalam hal ini, menelusuri bagian lain
bab ini -- Anda pasti akan memperhatikan berlimpahnya tanda $^{\ssbullet}$,
yang memberi halaman-halaman ini watak khas yang, saya duga, akan terasa
tidak sedap dipandang dan menakutkan, atau setidaknya melelahkan, bagi batin.
Banyak pengarang, mungkin sebagian besar, lebih suka menyederhanakan tipografi
dengan menggunakan simbol yang sama untuk fungsi dalam $\eusm L^0$ atau
$\eusm L^0_{\Sigma}$ dan untuk kelas ekuivalensinya dalam $L^0$; bahkan,
lazim pula digunakan sintaksis yang tidak membedakan keduanya, sehingga suatu
objek yang telah didefinisikan sebagai anggota $L^0$ tiba-tiba menjadi fungsi
dengan nilai-nilai aktual pada titik-titik ruang ukur dasarnya. Saya lebih
suka mempertahankan pembedaan yang tegas; Anda harus memutuskan sendiri apakah
hendak mengikuti saya. Karena saya telah memilih bentuk yang lebih panjang,
saya kira beban pembuktiannya ada pada saya untuk membenarkan keputusan
tersebut. (i) Semua orang akan setuju bahwa setidaknya terdapat perbedaan
formal antara sebuah fungsi dan sebuah himpunan fungsi. Hal ini sendiri tidak
membenarkan keharusan mempertahankan perbedaan tersebut dalam setiap kalimat;
uraian matematika akan mustahil jika kita selalu menuntut konsistensi dalam
pertanyaan semacam apakah (misalnya) bilangan $3$ yang termasuk dalam himpunan
$\Bbb N$ bilangan asli merupakan objek yang persis sama dengan bilangan $3$
yang termasuk dalam himpunan $\Bbb C$ bilangan kompleks, atau dengan ordinal
$3$. Namun, perbedaan antara sebuah objek dan sebuah himpunan yang memuatnya
merupakan perbedaan jenis yang cukup besar untuk membuat setiap kerancuan
sangat berbahaya; dan meskipun saya setuju bahwa Anda dapat mempelajari topik
ini tanpa menggunakan simbol berbeda untuk $f$ dan $f^{\ssbullet}$, saya
tidak berpendapat bahwa Anda dapat dengan aman menghindari pembedaan mental
itu lebih dari beberapa baris argumen. (ii) Sebagai pengajar, saya harus
mengatakan bahwa cukup banyak mahasiswa yang baru pertama kali menjumpai
bahan ini tersesat oleh kegagalan apa pun untuk membedakan $f$ dan
$f^{\ssbullet}$ hingga percaya bahwa pembedaan memang tidak perlu dibuat;
dan -- sebagai pengajar -- saya selalu meminta seorang mahasiswa meyakinkan
saya, dengan menuliskan bentuk argumen yang lebih teliti secara benar selama
beberapa minggu, bahwa ia memahami manipulasi yang diperlukan, sebelum saya
membiarkannya menempuh jalannya sendiri. (iii) Alasan mengapa pembedaan itu
{\it dapat} dihindari dalam jenis argumen tertentu ialah bahwa ruang Riesz
lengkap-$\sigma$ Dedekind $\eusm L^0_{\Sigma}$ sejajar sedemikian erat dengan
ruang Riesz lengkap-$\sigma$ Dedekind $L^0$ sehingga setiap proposisi yang
melibatkan hanya terhitung banyaknya anggota ruang-ruang ini kemungkinan
besar berlaku dalam yang satu jika dan hanya jika berlaku dalam yang lain.
Menurut pandangan saya, implikasi korespondensi ini berada tepat di jantung
teori ukuran. Karena itu, saya lebih suka membuatnya terus-menerus tampak
jelas, mengingatkan diri melalui simbolisme bahwa setiap teorema mempunyai
kembaran Siam, dan menyambut setiap tantangan untuk mengungkapkan teorema
kembar tersebut dalam bahasa yang sesuai. (iv) Ada cara-cara tertentu yang
membuat $\eusm L^0_{\Sigma}$ dan $L^0$ sungguh-sungguh sangat berbeda, dan
banyak gagasan menarik hanya dapat diungkapkan dalam bahasa yang memisahkan
keduanya dengan jelas.

Selama lebih dari separuh hidup saya hingga sekarang, saya merasa bahwa
butir-butir tersebut secara bersama-sama merupakan alasan yang cukup untuk
secara konsisten mempertahankan pembedaan formal antara $f$ dan
$f^{\ssbullet}$. Anda mungkin merasa bahwa dalam butir (iii) dan (iv) pada
paragraf terakhir saya mencoba mempertahankan dua posisi sekaligus; saya
berargumen bahwa baik kemiripan maupun perbedaan antara $L^0$ dan
$\eusm L^0$ mendukung pendirian saya. Memang, itulah persis pendirian saya.
Jika keduanya sama sekali berbeda, penggunaan bahasa yang sama untuk
keduanya tidak akan menimbulkan kerancuan; jika keduanya pada dasarnya sama,
tidak akan menjadi masalah jika kadang-kadang tidak jelas yang mana yang
sedang kita bicarakan.
}%end of notes

\discrpage
